\documentclass[10pt,a4paper]{amsart}
\usepackage[margin=19mm]{geometry}
\usepackage{amsmath,amssymb,mathtools}
\usepackage[scr=rsfs,cal=euler]{mathalfa}
\usepackage{xfrac}
\usepackage[unicode,hidelinks,bookmarks=false]{hyperref}
\newcommand{\ie}{i.e.\ }
\newcommand{\cf}{cf.\ }
\newcommand{\resp}{respectively\ }
\newcommand{\Subsection}{\subsection}
\providecommand{\nobreakdash}{\nobreak\hbox{-}}
\setlength{\emergencystretch}{2em}
\multlinegap0pt
\usepackage{tikz}
\usepackage{bm}
\usepackage{float}
\usepackage{amssymb}
\usepackage[shortlabels]{enumitem}
\usepackage{xcolor}
\newtheorem{lem}{Lemma}
\newtheorem{prop}{Proposition}
\newcommand{\F}{\mathbb{F}}
\DeclareMathOperator{\Mat}{Mat}
\newcommand{\fp}{\mathfrak{p}}
\title{Rank-metric codes from Drinfeld modules}
\author{Giacomo Micheli, Mihran Papikian}
\date{Source: arXiv:2601.03653v2 (CC BY 4.0)}
\begin{document}
\maketitle

\noindent\emph{Source and licence.} Rank-metric codes from Drinfeld modules. Version arXiv:2601.03653v2 (CC BY 4.0), distributed by arXiv under the Creative Commons Attribution 4.0 International licence, http://creativecommons.org/licenses/by/4.0/, as recorded in the arXiv licence metadata for this exact version and shown on the abstract page https://arxiv.org/abs/2601.03653v2. Full licence notice: \texttt{licenses/arxiv-2601.03653-CC-BY-4.0.txt}.\par\smallskip

\noindent\textit{Editorial scope.} Counted unit: the article's prop propCM (source0.tex lines 1430--1436). The article's complete original proof of it is delivered with it (source0.tex lines 1437--1471, one \texttt{proof} environment of 35 lines). No mathematics is written, repaired or re-derived: the statement and the proof are the article's own lines, in the article's own notation, and no retained line is substituted or re-flowed.  Retained uncounted as the source-local context the proof needs: lines 1418--1423.  Nothing was omitted from any retained span, so the package carries no omission record.  No retained line contains a citation, so the wrapper carries no bibliography and no bibliography entry.  No retained line contains a figure environment, an image-inclusion command or drawing code, so no image is delivered and the package declares no asset dependency.  Editorial formatting only, in addition: an AMS \texttt{amsart} preamble that restates the article's own declarations and macros used by the retained lines, namely the environments \texttt{lem}, \texttt{prop} on separate counters, together with the standard packages those lines need.  The retained environments print in this document as Lemma 1, Proposition 1 in order of appearance, whereas the article prints the same items with the numbering of its own full document.  The complete native source of the article as distributed in the arXiv e-print for this exact version is bundled unchanged as \texttt{sources/192238/source0.tex}.
 \begin{lem}\label{lemC(F_p)}
 	The image of $\iota$ satisfies 
 	\[
 	\iota(\Mat_r(\F_\fp)) = C_{\Mat_{rd}(\F_q)}(\iota(\F_\fp)). 
 	\]
 \end{lem}

 \begin{prop}\label{propCM}
 	Let $M\in \Mat_r(\F_\fp)$. The following are equivalent: 
 	\begin{itemize}
 		\item[(i)] $C_{\Mat_{rd}(\F_q)}(\iota(M))=\iota\left(C_{\Mat_r(\F_\fp)}(M)\right)$.
 		\item[(ii)]  $\iota(\F_\fp)\subseteq \F_q[\iota(M)]$. 
 	\end{itemize}
 \end{prop}

 \begin{proof} To simplify the notation, we will denote 
 	\[C(\iota(M))= C_{\Mat_{rd}(\F_q)}(\iota(M))\quad \text{and}\quad \iota(C(M))=\iota\left(C_{\Mat_r(\F_\fp)}(M)\right).
 	\] 
 	
 	If $N\in \Mat_r(\F_\fp)$ satisfies $NM=MN$, then applying $\iota$ (which is a ring homomorphism) gives $\iota(N)\iota(M)=\iota(M)\iota(N)$. 
 	Thus, the inclusion  
 	\[
 	\iota\left(C(M)\right) \subseteq C(\iota(M)) 
 	\]
 	holds unconditionally. 
 	
 	For the reverse inclusion, by Lemma \ref{lemC(F_p)}, we have 
 	\[
 	\iota\left(C(M)\right)  = 	\iota(\Mat_r(\F_\fp))  \cap C(\iota(M)) = C(\iota(\F_\fp)) \cap C(\iota(M)). 
 	\]
 	Therefore, the equality $\iota\left(C(M)\right)  = C(\iota(M))$ holds if and only if 
 	\[
 	C(\iota(M))\subseteq C(\iota(\F_\fp)). 
 	\]
 	
 	(ii)$\Rightarrow$(i): Suppose $\iota(\F_\fp)\subseteq \F_q[\iota(M)]$. If $X\in C(\iota(M))$, then $X$ commutes with every 
 	polynomial in $\iota(M)$, hence $X$ commutes with every element of $\F_q[\iota(M)]$. In particular, $X$ commutes 
 	with every element of $\iota(\F_\fp)$, so $X\in C(\iota(\F_\fp))$. 
 	
 	(i)$\Rightarrow$(ii):  Suppose $C(\iota(M))\subseteq C(\iota(\F_\fp))$. Taking centralizers reverses inclusions, so 
 	\[
 	C(C(\iota(\F_\fp)))\subseteq C(C(\iota(M))). 
 	\]
 	On the other hand, by the Double Centralizer Theorem, 
 	\begin{align*}
 	C(C(\iota(\F_\fp))) &= \iota(\F_\fp), \\ 
 	C(C(\iota(M))) &= \F_q[\iota(M)].  
 	\end{align*}
 	Therefore, $\iota(\F_\fp) \subseteq \F_q[\iota(M)]$. 
 \end{proof}

\end{document}
